%%%%%%%%%%%%%%%%%%%%%%%%%%%%%%%%%%%%%%%%%%%%%%%%%%%%%%%%%%%%%
%% CATEGORY THEORY — Preamble + Chapters 1–2
%% Categories, Functors, Natural Transformations / Duality
%%%%%%%%%%%%%%%%%%%%%%%%%%%%%%%%%%%%%%%%%%%%%%%%%%%%%%%%%%%%%

\documentclass[a4paper,12pt]{book}

%% ---- Encoding and language ----
\usepackage[utf8]{inputenc}
\usepackage[T1]{fontenc}
\usepackage[english]{babel}

%% ---- Mathematics ----
\usepackage{amsmath}
\usepackage{amssymb}
\usepackage{amsthm}
\usepackage{mathtools}
\usepackage{stmaryrd}

%% ---- Coloured boxes ----
\usepackage[theorems,skins,breakable]{tcolorbox}

%% ---- Diagrams and graphics ----
\usepackage{tikz}
\usepackage{tikz-cd}
\usepackage{pgfplots}
\pgfplotsset{compat=1.18}
\usepackage{graphicx}
\usepackage{float}
\usepackage{caption}

%% ---- Lists, tables ----
\usepackage{enumitem}
\usepackage{booktabs}
\usepackage{array}
\usepackage{longtable}

%% ---- Cross-references and indexing ----
\usepackage{hyperref}
\hypersetup{
    colorlinks=true,
    linkcolor=blue!70!black,
    citecolor=green!50!black,
    urlcolor=blue!60!black
}
\usepackage{cleveref}
\usepackage{makeidx}
\usepackage{minitoc}

\makeindex

%% ============================================================
%%  Math operators
%% ============================================================
\DeclareMathOperator{\Hom}{Hom}
\DeclareMathOperator{\End}{End}
\DeclareMathOperator{\Aut}{Aut}
\DeclareMathOperator{\im}{im}
\DeclareMathOperator{\Id}{Id}
\DeclareMathOperator{\Ker}{Ker}
\DeclareMathOperator{\coker}{coker}
\DeclareMathOperator{\Ob}{Ob}
\DeclareMathOperator{\Mor}{Mor}
\DeclareMathOperator{\Fun}{Fun}
\DeclareMathOperator{\Nat}{Nat}
\DeclareMathOperator{\colim}{colim}
\DeclareMathOperator{\eq}{eq}
\DeclareMathOperator{\coeq}{coeq}

%% ---- Named categories ----
\newcommand{\Set}{\mathbf{Set}}
\newcommand{\Grp}{\mathbf{Grp}}
\newcommand{\Ab}{\mathbf{Ab}}
\newcommand{\Ring}{\mathbf{Ring}}
\newcommand{\Top}{\mathbf{Top}}
\newcommand{\Vect}{\mathbf{Vect}}
\newcommand{\Mod}{\mathbf{Mod}}
\newcommand{\Cat}{\mathbf{Cat}}

%% ---- Convenience commands ----
\newcommand{\op}{\mathrm{op}}
\newcommand{\id}{\mathrm{id}}
\newcommand{\N}{\mathbb{N}}
\newcommand{\Z}{\mathbb{Z}}
\newcommand{\Q}{\mathbb{Q}}
\newcommand{\R}{\mathbb{R}}
\newcommand{\C}{\mathbb{C}}
\newcommand{\K}{\mathbb{K}}

%% ============================================================
%%  Theorem environments (amsthm + tcolorbox styling)
%% ============================================================

%% --- amsthm declarations ---
\theoremstyle{definition}
\newtheorem{definition}{Definition}[section]
\newtheorem{example}[definition]{Example}
\newtheorem{exercise}[definition]{Exercise}

\theoremstyle{plain}
\newtheorem{theorem}[definition]{Theorem}
\newtheorem{proposition}[definition]{Proposition}
\newtheorem{lemma}[definition]{Lemma}
\newtheorem{corollary}[definition]{Corollary}

\theoremstyle{remark}
\newtheorem{remark}[definition]{Remark}
\newtheorem{notation}[definition]{Notation}

%% --- tcolorbox styling ---
\tcolorboxenvironment{definition}{
    enhanced, breakable,
    colback=blue!3, colframe=blue!50!black,
    boxrule=0.6pt, arc=2pt,
    left=6pt, right=6pt, top=4pt, bottom=4pt,
    fonttitle=\bfseries
}
\tcolorboxenvironment{example}{
    enhanced, breakable,
    colback=green!3, colframe=green!40!black,
    boxrule=0.5pt, arc=2pt,
    left=6pt, right=6pt, top=4pt, bottom=4pt,
    fonttitle=\bfseries
}
\tcolorboxenvironment{exercise}{
    enhanced, breakable,
    colback=orange!4, colframe=orange!60!black,
    boxrule=0.5pt, arc=2pt,
    left=6pt, right=6pt, top=4pt, bottom=4pt,
    fonttitle=\bfseries
}
\tcolorboxenvironment{theorem}{
    enhanced, breakable,
    colback=red!3, colframe=red!60!black,
    boxrule=0.7pt, arc=2pt,
    left=6pt, right=6pt, top=4pt, bottom=4pt,
    fonttitle=\bfseries
}
\tcolorboxenvironment{proposition}{
    enhanced, breakable,
    colback=red!2, colframe=red!45!black,
    boxrule=0.6pt, arc=2pt,
    left=6pt, right=6pt, top=4pt, bottom=4pt,
    fonttitle=\bfseries
}
\tcolorboxenvironment{lemma}{
    enhanced, breakable,
    colback=yellow!4, colframe=yellow!50!black,
    boxrule=0.5pt, arc=2pt,
    left=6pt, right=6pt, top=4pt, bottom=4pt,
    fonttitle=\bfseries
}
\tcolorboxenvironment{corollary}{
    enhanced, breakable,
    colback=purple!3, colframe=purple!40!black,
    boxrule=0.5pt, arc=2pt,
    left=6pt, right=6pt, top=4pt, bottom=4pt,
    fonttitle=\bfseries
}
\tcolorboxenvironment{remark}{
    enhanced, breakable,
    colback=gray!4, colframe=gray!50!black,
    boxrule=0.4pt, arc=2pt,
    left=6pt, right=6pt, top=4pt, bottom=4pt,
    fonttitle=\bfseries
}
\tcolorboxenvironment{notation}{
    enhanced, breakable,
    colback=cyan!3, colframe=cyan!40!black,
    boxrule=0.4pt, arc=2pt,
    left=6pt, right=6pt, top=4pt, bottom=4pt,
    fonttitle=\bfseries
}

%% ============================================================
%%  Title
%% ============================================================
\title{\Huge\textbf{Category Theory}\\[0.5em]
       \Large A Graduate Course}
\author{}
\date{}

%% ============================================================
\begin{document}
%% ============================================================

\maketitle
\dominitoc
\tableofcontents

%% ============================================================
%%  PREFACE
%% ============================================================
\chapter*{Preface}\label{ch:preface}
\addcontentsline{toc}{chapter}{Preface}

Category theory provides a unifying language for mathematics.  Originally
introduced by Eilenberg and Mac\,Lane in the 1940s to formalise natural
transformations in algebraic topology, it has since become an
indispensable tool in algebra, geometry, logic, computer science, and
mathematical physics.

This text is designed for a one-semester graduate course.  We assume
familiarity with basic algebra (groups, rings, modules) and point-set
topology, but no prior knowledge of category theory itself.  The
exposition emphasises both the abstract framework and its concrete
manifestations: every definition is illustrated by examples drawn from
algebra, topology, and analysis.

\medskip
\noindent\textbf{Organisation.}
Chapters~1 and~2 lay the foundations: categories, functors, natural
transformations, and the duality principle.
Chapters~3 and~4 develop limits, colimits, and adjunctions.
Later chapters treat representability, the Yoneda lemma, Kan extensions,
abelian categories, and monoidal categories.

\medskip
\noindent\textbf{Exercises.}
Each chapter ends with a graded exercise set.  Working through these
is essential for mastering the material.

%% ============================================================
%%  NOTATION
%% ============================================================
\chapter*{Notation and Conventions}\label{ch:notation}
\addcontentsline{toc}{chapter}{Notation and Conventions}

\begin{longtable}{>{\raggedright}p{4cm} p{9.5cm}}
\toprule
\textbf{Symbol} & \textbf{Meaning} \\
\midrule
\endhead
$\N,\;\Z,\;\Q,\;\R,\;\C$ &
    Natural numbers (including~$0$), integers, rationals, reals,
    complex numbers \\[4pt]
$\K$ &
    An arbitrary field \\[4pt]
$\Set,\;\Grp,\;\Ab$ &
    Categories of sets, groups, abelian groups \\[4pt]
$\Ring,\;\Top,\;\Vect_\K$ &
    Categories of rings, topological spaces,
    $\K$-vector spaces \\[4pt]
$\Mod_R$ &
    Category of left $R$-modules \\[4pt]
$\Cat$ &
    Category of small categories \\[4pt]
$\Hom_\mathcal{C}(A,B)$ &
    Set of morphisms $A\to B$ in~$\mathcal{C}$ \\[4pt]
$\mathcal{C}^{\op}$ &
    Opposite (dual) category of~$\mathcal{C}$ \\[4pt]
$F\colon\mathcal{C}\to\mathcal{D}$ &
    Functor from $\mathcal{C}$ to $\mathcal{D}$ \\[4pt]
$\alpha\colon F\Rightarrow G$ &
    Natural transformation from $F$ to $G$ \\[4pt]
$\id_A$ &
    Identity morphism on~$A$ \\[4pt]
$g\circ f$ &
    Composition: first~$f$, then~$g$ \\[4pt]
$[\mathcal{C},\mathcal{D}]$ or $\Fun(\mathcal{C},\mathcal{D})$ &
    Functor category \\[4pt]
\bottomrule
\end{longtable}

\noindent
Throughout, ``category'' means a locally small category unless otherwise
stated.  We use the terms ``morphism'', ``arrow'', and ``map''
interchangeably.  All rings are assumed to have a multiplicative
identity.


%%%%%%%%%%%%%%%%%%%%%%%%%%%%%%%%%%%%%%%%%%%%%%%%%%%%%%%%%%%%
%%       CHAPTER 1: CATEGORIES, FUNCTORS, NATURAL         %%
%%                  TRANSFORMATIONS                        %%
%%%%%%%%%%%%%%%%%%%%%%%%%%%%%%%%%%%%%%%%%%%%%%%%%%%%%%%%%%%%

\chapter{Categories, Functors, and Natural Transformations}
\label{ch:cat-fun-nat}
\minitoc

\vspace{1em}

Categories, functors, and natural transformations constitute the three
fundamental notions of category theory.  A category axiomatises the idea
of ``objects with structure-preserving maps between them''.  A functor is
a morphism of categories, and a natural transformation is a morphism of
functors.  Already at this first level of abstraction one finds a
striking pattern: the passage from objects to morphisms between objects
repeats at every stage.

In this chapter we introduce these three notions carefully, give
copious examples, and establish basic terminology---monomorphisms,
epimorphisms, isomorphisms, initial and terminal objects---that will be
used throughout the course.

\index{category|(}


%% ============================================================
\section{Categories}\label{sec:categories}
%% ============================================================

\begin{definition}[Category]\label{def:category}
\index{category!definition}
\index{object}
\index{morphism}
\index{composition}
\index{identity morphism}
A \textbf{category}~$\mathcal{C}$ consists of the following data:
\begin{enumerate}[label=(\roman*)]
    \item A collection $\Ob(\mathcal{C})$ of \emph{objects}.
    \item For every ordered pair of objects $A,B\in\Ob(\mathcal{C})$,
          a set $\Hom_\mathcal{C}(A,B)$ of \emph{morphisms}
          (or \emph{arrows}) from~$A$ to~$B$.
          We write $f\colon A\to B$ to indicate
          $f\in\Hom_\mathcal{C}(A,B)$.
    \item For every triple of objects $A,B,C$, a \emph{composition law}
          \[
              \circ\colon \Hom_\mathcal{C}(B,C)\times
                          \Hom_\mathcal{C}(A,B)
                          \longrightarrow \Hom_\mathcal{C}(A,C),
              \qquad (g,f)\longmapsto g\circ f.
          \]
    \item For every object $A$, an \emph{identity morphism}
          $\id_A\in\Hom_\mathcal{C}(A,A)$.
\end{enumerate}
These data are subject to two axioms:
\begin{enumerate}[label=\textbf{(C\arabic*)}]
    \item \textbf{Associativity.}\index{associativity}
          For all $f\colon A\to B$, $g\colon B\to C$,
          $h\colon C\to D$,
          \[
              h\circ(g\circ f) = (h\circ g)\circ f.
          \]
    \item \textbf{Identity.}
          For all $f\colon A\to B$,
          \[
              f\circ\id_A = f = \id_B\circ f.
          \]
\end{enumerate}
\end{definition}

\begin{remark}\label{rem:hom-disjoint}
\index{hom-set}
We require the hom-sets to be pairwise disjoint: every morphism~$f$ has a
uniquely determined \emph{domain} (or \emph{source}) and \emph{codomain}
(or \emph{target}).  Some authors encode this by defining a morphism as a
triple $(f,A,B)$.
\end{remark}

\begin{notation}\label{not:hom-notation}
\index{hom-set!notation}
We write $\Hom_\mathcal{C}(A,B)$, or $\mathcal{C}(A,B)$, or
$\Mor_\mathcal{C}(A,B)$ interchangeably.
When the ambient category is clear we simply write $\Hom(A,B)$.
\end{notation}


%% ============================================================
\section{First examples of categories}\label{sec:first-examples}
%% ============================================================

\begin{example}[Algebraic categories]\label{ex:algebraic-cats}
\index{category!Set@$\Set$}
\index{category!Grp@$\Grp$}
\index{category!Ab@$\Ab$}
\index{category!Ring@$\Ring$}
\index{category!Mod@$\Mod_R$}
\index{category!Vect@$\Vect_\K$}
The following are categories:
\begin{enumerate}[label=(\roman*)]
    \item $\Set$: objects are sets, morphisms are functions.
    \item $\Grp$: objects are groups, morphisms are group homomorphisms.
    \item $\Ab$: objects are abelian groups, morphisms are group
          homomorphisms.
    \item $\Ring$: objects are (unital) rings, morphisms are ring
          homomorphisms (preserving~$1$).
    \item $\Mod_R$: for a ring~$R$, objects are left $R$-modules,
          morphisms are $R$-linear maps.
    \item $\Vect_\K$: objects are $\K$-vector spaces, morphisms are
          $\K$-linear maps.  This is the special case
          $\Mod_\K$.
\end{enumerate}
In each case, composition is the usual composition of functions and the
identity morphism on an object~$A$ is the identity function
$\id_A\colon A\to A$.
\end{example}

\begin{example}[Topological category]\label{ex:top-cat}
\index{category!Top@$\Top$}
$\Top$: objects are topological spaces, morphisms are continuous maps.
\end{example}

\begin{example}[Ordered sets as categories]\label{ex:poset-cat}
\index{category!poset}
\index{poset!as category}
\index{preorder!as category}
Let $(P,\le)$ be a preorder (a set equipped with a reflexive, transitive
relation).  Define a category~$\mathcal{C}_P$ by:
\begin{itemize}
    \item $\Ob(\mathcal{C}_P) = P$.
    \item For $a,b\in P$:
    \[
        \Hom(a,b) =
        \begin{cases}
            \{\ast\} & \text{if } a\le b, \\
            \varnothing & \text{otherwise.}
        \end{cases}
    \]
\end{itemize}
Transitivity gives composition; reflexivity gives identities.
If $(P,\le)$ is a partial order, distinct objects are never isomorphic.
\end{example}

\begin{example}[Monoids as categories]\label{ex:monoid-cat}
\index{monoid!as category}
\index{category!one-object}
A monoid $(M,\cdot,e)$ may be viewed as a category with a single
object~$\ast$ and $\Hom(\ast,\ast)=M$.
Composition is the monoid operation and the identity morphism is~$e$.

Conversely, every category with exactly one object is a monoid.
A group is a one-object category in which every morphism is invertible.
\end{example}

\begin{example}[The category $\Cat$]\label{ex:cat-of-cats}
\index{category!Cat@$\Cat$}
The category $\Cat$ has small categories as objects and functors
(defined in 
ef{sec:functors}) as morphisms.
\end{example}


%% ============================================================
\section{Small and locally small categories}\label{sec:size}
%% ============================================================

\begin{definition}[Size conditions]\label{def:size}
\index{small category}
\index{locally small category}
\index{large category}
A category~$\mathcal{C}$ is called:
\begin{enumerate}[label=(\roman*)]
    \item \textbf{small} if $\Ob(\mathcal{C})$ is a set
          (not a proper class);
    \item \textbf{locally small} if for every pair of objects
          $A,B$, the collection $\Hom_\mathcal{C}(A,B)$ is a set.
\end{enumerate}
A category that is not small is called \textbf{large}.
\end{definition}

\begin{remark}\label{rem:size-remarks}
Every small category is locally small.
The categories $\Set$, $\Grp$, $\Top$, etc.\ are locally small but not
small (their collection of objects is a proper class).
Every preorder viewed as a category is small provided the underlying
set is a set.
Throughout this text, ``category'' means ``locally small category''
unless otherwise stated.
\end{remark}


%% ============================================================
\section{Monomorphisms, epimorphisms, and isomorphisms}
\label{sec:mono-epi-iso}
%% ============================================================

\begin{definition}[Monomorphism]\label{def:mono}
\index{monomorphism}
\index{monic|see{monomorphism}}
A morphism $f\colon A\to B$ in a category~$\mathcal{C}$ is a
\textbf{monomorphism} (or is \emph{monic}) if for every pair of
morphisms $g_1,g_2\colon C\to A$,
\[
    f\circ g_1 = f\circ g_2 \implies g_1 = g_2.
\]
We denote a monomorphism by $f\colon A\rightarrowtail B$.
\end{definition}

\begin{definition}[Epimorphism]\label{def:epi}
\index{epimorphism}
\index{epic|see{epimorphism}}
A morphism $f\colon A\to B$ is an \textbf{epimorphism} (or is
\emph{epic}) if for every pair $h_1,h_2\colon B\to D$,
\[
    h_1\circ f = h_2\circ f \implies h_1 = h_2.
\]
We denote an epimorphism by $f\colon A\twoheadrightarrow B$.
\end{definition}

\begin{definition}[Isomorphism]\label{def:iso}
\index{isomorphism}
\index{inverse morphism}
A morphism $f\colon A\to B$ is an \textbf{isomorphism} if there exists a
morphism $g\colon B\to A$ such that
\[
    g\circ f = \id_A \qquad\text{and}\qquad f\circ g = \id_B.
\]
The morphism~$g$ is unique and is called the \emph{inverse} of~$f$,
written $f^{-1}$.  Two objects are \emph{isomorphic},
$A\cong B$\index{isomorphic objects}, if there exists an isomorphism
between them.
\end{definition}

\begin{proposition}\label{prop:iso-mono-epi}
\index{isomorphism!versus monomorphism and epimorphism}
Every isomorphism is both a monomorphism and an epimorphism.
\end{proposition}

\begin{proof}
Let $f\colon A\to B$ be an isomorphism with inverse~$g$.
If $f\circ g_1 = f\circ g_2$, apply~$g$ on the left:
$g_1 = g\circ f\circ g_1 = g\circ f\circ g_2 = g_2$,
so~$f$ is monic.  Dually, $f$ is epic.
\end{proof}

\begin{remark}\label{rem:mono-epi-not-iso}
\index{epimorphism!not surjective}
The converse fails in general.  In~$\Ring$, the inclusion
$\Z\hookrightarrow\Q$ is both monic and epic, but it is not an
isomorphism.
\end{remark}

\begin{example}[Mono and epi in concrete categories]\label{ex:mono-epi-concrete}
\index{monomorphism!in Set@in $\Set$}
\index{epimorphism!in Set@in $\Set$}
\begin{enumerate}[label=(\roman*)]
    \item In $\Set$: monomorphisms are injections, epimorphisms are
          surjections, isomorphisms are bijections.
    \item In $\Grp$: monomorphisms are injective homomorphisms,
          epimorphisms are surjective homomorphisms.
    \item In $\Top$: monomorphisms are injective continuous maps;
          epimorphisms are surjective continuous maps
          (but \emph{not} necessarily quotient maps).
\end{enumerate}
\end{example}

\begin{lemma}\label{lem:mono-epi-composition}
\index{monomorphism!composition}
\index{epimorphism!composition}
Let $f\colon A\to B$ and $g\colon B\to C$ be morphisms.
\begin{enumerate}[label=(\roman*)]
    \item If $g\circ f$ is monic, then $f$ is monic.
    \item If $g\circ f$ is epic, then $g$ is epic.
    \item If both $f$ and $g$ are monic (resp.\ epic), then so is
          $g\circ f$.
\end{enumerate}
\end{lemma}

\begin{proof}
(i) Suppose $g\circ f$ is monic and $f\circ h_1 = f\circ h_2$.
Then $g\circ f\circ h_1 = g\circ f\circ h_2$, hence $h_1=h_2$.
Parts~(ii) and~(iii) are similar.
\end{proof}


%% ============================================================
\section{Initial and terminal objects}\label{sec:initial-terminal}
%% ============================================================

\begin{definition}[Initial object]\label{def:initial}
\index{initial object}
\index{object!initial}
An object $I\in\mathcal{C}$ is \textbf{initial} if for every object
$A\in\mathcal{C}$ there exists a unique morphism $I\to A$.
\end{definition}

\begin{definition}[Terminal object]\label{def:terminal}
\index{terminal object}
\index{object!terminal}
An object $T\in\mathcal{C}$ is \textbf{terminal} if for every object
$A\in\mathcal{C}$ there exists a unique morphism $A\to T$.
\end{definition}

\begin{definition}[Zero object]\label{def:zero-object}
\index{zero object}
An object that is both initial and terminal is called a
\textbf{zero object}.
\end{definition}

\begin{proposition}\label{prop:initial-terminal-unique}
\index{initial object!uniqueness}
\index{terminal object!uniqueness}
If an initial (resp.\ terminal, resp.\ zero) object exists, it is unique
up to unique isomorphism.
\end{proposition}

\begin{proof}
Let $I$ and $I'$ be initial.  By the universal property there exist
unique morphisms $f\colon I\to I'$ and $g\colon I'\to I$.
Then $g\circ f\colon I\to I$ must equal $\id_I$ (uniqueness of the
morphism $I\to I$), and similarly $f\circ g = \id_{I'}$.
The argument for terminal objects is dual.
\end{proof}

\begin{example}[Initial and terminal objects]\label{ex:initial-terminal}
\index{empty set!as initial object}
\begin{enumerate}[label=(\roman*)]
    \item In $\Set$: the empty set $\varnothing$ is initial; any
          singleton $\{*\}$ is terminal.
    \item In $\Grp$ and $\Ab$: the trivial group $\{e\}$ is a
          zero object.
    \item In $\Ring$: $\Z$ is initial (unique ring homomorphism
          $\Z\to R$ sending $1\mapsto 1_R$); the zero ring~$\{0\}$
          is terminal.
    \item In $\Top$: the empty space is initial; any one-point space
          is terminal.
    \item In $\Vect_\K$: the zero space $\{0\}$ is a zero object.
\end{enumerate}
\end{example}

\begin{remark}\label{rem:zero-morphism}
\index{zero morphism}
If $\mathcal{C}$ has a zero object~$0$, then for any objects $A,B$ there
is a unique morphism $A\to 0\to B$, called the \emph{zero morphism}
and denoted $0_{AB}$ (or simply~$0$).
\end{remark}


%% ============================================================
\section{Functors}\label{sec:functors}
%% ============================================================

\index{functor|(}

\begin{definition}[Functor]\label{def:functor}
\index{functor!covariant}
\index{covariant functor|see{functor}}
Let $\mathcal{C}$ and $\mathcal{D}$ be categories.
A \textbf{(covariant) functor} $F\colon\mathcal{C}\to\mathcal{D}$
consists of:
\begin{enumerate}[label=(\roman*)]
    \item A map on objects:
          $A\mapsto F(A)$ for each $A\in\Ob(\mathcal{C})$.
    \item A map on morphisms: for each $f\colon A\to B$ in~$\mathcal{C}$,
          a morphism $F(f)\colon F(A)\to F(B)$ in~$\mathcal{D}$.
\end{enumerate}
These must satisfy:
\begin{enumerate}[label=\textbf{(F\arabic*)}]
    \item $F(\id_A) = \id_{F(A)}$ for every object $A$.
    \item $F(g\circ f) = F(g)\circ F(f)$ for every composable pair
          $f,g$.
\end{enumerate}
\end{definition}

\begin{definition}[Contravariant functor]\label{def:contravariant}
\index{functor!contravariant}
\index{contravariant functor}
A \textbf{contravariant functor} $F\colon\mathcal{C}\to\mathcal{D}$ is
a (covariant) functor $F\colon\mathcal{C}^{\op}\to\mathcal{D}$.
Equivalently, $F$ reverses the direction of morphisms:
$F(f)\colon F(B)\to F(A)$ when $f\colon A\to B$, and
$F(g\circ f) = F(f)\circ F(g)$.
\end{definition}

\begin{example}[Forgetful functors]\label{ex:forgetful}
\index{functor!forgetful}
\index{forgetful functor}
\begin{enumerate}[label=(\roman*)]
    \item $U\colon\Grp\to\Set$: sends a group to its underlying set
          and a homomorphism to the underlying function.
    \item $U\colon\Top\to\Set$: forgets the topology.
    \item $U\colon\Vect_\K\to\Ab$: forgets scalar multiplication,
          retaining only the additive group.
    \item $U\colon\Ring\to\Ab$: sends a ring to its underlying
          additive group.
\end{enumerate}
\end{example}

\begin{example}[Free functors]\label{ex:free-functor}
\index{functor!free}
\index{free functor}
\begin{enumerate}[label=(\roman*)]
    \item $F\colon\Set\to\Grp$: sends a set~$S$ to the free group on~$S$.
    \item $F\colon\Set\to\Vect_\K$: sends a set~$S$ to the vector
          space with basis~$S$.
    \item $F\colon\Set\to\Ab$: sends a set~$S$ to the free abelian
          group $\Z^{(S)} = \bigoplus_{s\in S}\Z$.
\end{enumerate}
\end{example}

\begin{example}[Power-set functor]\label{ex:powerset-functor}
\index{functor!power-set}
\index{power-set functor}
There are two natural functors $\Set\to\Set$ associated with the power
set:
\begin{enumerate}[label=(\roman*)]
    \item \textbf{Covariant:} $\mathcal{P}\colon\Set\to\Set$ sends a
          set~$S$ to $\mathcal{P}(S)$ and a function $f\colon S\to T$
          to the direct image $f_*\colon \mathcal{P}(S)\to\mathcal{P}(T)$,
          $A\mapsto f(A)$.
    \item \textbf{Contravariant:} $\mathcal{P}^{\op}\colon\Set^{\op}\to\Set$
          sends $f\colon S\to T$ to the pre-image
          $f^*\colon\mathcal{P}(T)\to\mathcal{P}(S)$,
          $B\mapsto f^{-1}(B)$.
\end{enumerate}
\end{example}

\begin{example}[Hom-functors]\label{ex:hom-functors}
\index{functor!hom}
\index{hom-functor}
For any locally small category~$\mathcal{C}$ and object $A\in\mathcal{C}$:
\begin{enumerate}[label=(\roman*)]
    \item The \textbf{covariant hom-functor}
          $\Hom_\mathcal{C}(A,-)\colon\mathcal{C}\to\Set$
          sends $B\mapsto\Hom(A,B)$ and $f\colon B\to C$ to
          $f_*\colon\Hom(A,B)\to\Hom(A,C)$, $g\mapsto f\circ g$.
    \item The \textbf{contravariant hom-functor}
          $\Hom_\mathcal{C}(-,B)\colon\mathcal{C}^{\op}\to\Set$
          sends $A\mapsto\Hom(A,B)$ and $f\colon A'\to A$ to
          $f^*\colon\Hom(A,B)\to\Hom(A',B)$, $g\mapsto g\circ f$.
\end{enumerate}
\end{example}

\begin{example}[Fundamental group functor]\label{ex:pi1-functor}
\index{functor!fundamental group}
\index{fundamental group!functor}
Let $\Top_*$ denote the category of pointed topological spaces and
base-point-preserving continuous maps.  The fundamental group construction
defines a functor
\[
    \pi_1\colon\Top_*\longrightarrow\Grp,
    \qquad (X,x_0)\longmapsto\pi_1(X,x_0).
\]
A continuous map $f\colon(X,x_0)\to(Y,y_0)$ induces the group
homomorphism $f_*\colon\pi_1(X,x_0)\to\pi_1(Y,y_0)$ given by
$[\gamma]\mapsto[f\circ\gamma]$.
\end{example}


%% ============================================================
\section{Faithful, full, and essentially surjective functors}
\label{sec:faithful-full}
%% ============================================================

\begin{definition}[Faithful, full, fully faithful]\label{def:faithful-full}
\index{functor!faithful}
\index{functor!full}
\index{functor!fully faithful}
\index{faithful functor}
\index{full functor}
A functor $F\colon\mathcal{C}\to\mathcal{D}$ is called:
\begin{enumerate}[label=(\roman*)]
    \item \textbf{faithful} if for every pair $A,B\in\mathcal{C}$
          the map
          \[
              F_{A,B}\colon\Hom_\mathcal{C}(A,B)
              \longrightarrow\Hom_\mathcal{D}(F(A),F(B))
          \]
          is injective;
    \item \textbf{full} if every $F_{A,B}$ is surjective;
    \item \textbf{fully faithful} if every $F_{A,B}$ is bijective.
\end{enumerate}
\end{definition}

\begin{definition}[Essentially surjective]\label{def:ess-surj}
\index{functor!essentially surjective}
\index{essentially surjective}
A functor $F\colon\mathcal{C}\to\mathcal{D}$ is
\textbf{essentially surjective} (or \emph{essentially surjective on
objects}) if for every $D\in\mathcal{D}$ there exists $C\in\mathcal{C}$
with $F(C)\cong D$.
\end{definition}

\begin{example}\label{ex:faithful-full-examples}
\begin{enumerate}[label=(\roman*)]
    \item Every forgetful functor (e.g.\ $U\colon\Grp\to\Set$) is
          faithful.  It is not full in general: not every function
          between the underlying sets of two groups is a homomorphism.
    \item The inclusion functor $\Ab\hookrightarrow\Grp$ is fully
          faithful: a homomorphism between abelian groups is the same
          thing whether viewed in~$\Ab$ or~$\Grp$.
    \item A functor is an \emph{equivalence of categories}
          if and only if it is fully faithful and essentially surjective
          (assuming the axiom of choice).
\end{enumerate}
\end{example}

\begin{proposition}\label{prop:ff-reflects-iso}
\index{functor!fully faithful!reflects isomorphisms}
A fully faithful functor reflects isomorphisms: if
$F(f)$ is an isomorphism in~$\mathcal{D}$, then $f$ is an isomorphism
in~$\mathcal{C}$.
\end{proposition}

\begin{proof}
Let $F(f)\colon F(A)\to F(B)$ be an isomorphism with inverse~$h$.
Since $F$ is full, $h = F(g)$ for some $g\colon B\to A$.
Then $F(g\circ f) = F(g)\circ F(f) = \id_{F(A)} = F(\id_A)$.
Since $F$ is faithful, $g\circ f = \id_A$.  Similarly
$f\circ g = \id_B$.
\end{proof}


%% ============================================================
\section{Natural transformations}\label{sec:nat-trans}
%% ============================================================

\index{natural transformation|(}

\begin{definition}[Natural transformation]\label{def:nat-trans}
\index{natural transformation!definition}
\index{component!of a natural transformation}
Let $F,G\colon\mathcal{C}\to\mathcal{D}$ be functors.
A \textbf{natural transformation} $\alpha\colon F\Rightarrow G$ is a
family of morphisms in~$\mathcal{D}$,
\[
    \bigl\{\alpha_A\colon F(A)\to G(A)\bigr\}_{A\in\Ob(\mathcal{C})},
\]
called the \emph{components} of~$\alpha$, such that for every morphism
$f\colon A\to B$ in~$\mathcal{C}$ the following \emph{naturality square}
commutes:
\[
\begin{tikzcd}
    F(A) \arrow[r, "\alpha_A"] \arrow[d, "F(f)"'] &
    G(A) \arrow[d, "G(f)"] \\
    F(B) \arrow[r, "\alpha_B"'] &
    G(B)
\end{tikzcd}
\]
That is, $G(f)\circ\alpha_A = \alpha_B\circ F(f)$ for all
$f\colon A\to B$.
\end{definition}

\begin{definition}[Natural isomorphism]\label{def:nat-iso}
\index{natural isomorphism}
\index{natural transformation!isomorphism}
A natural transformation $\alpha\colon F\Rightarrow G$ is a
\textbf{natural isomorphism} if every component $\alpha_A$ is an
isomorphism.  We then write $F\cong G$ and say $F$ and $G$ are
\emph{naturally isomorphic}.
\end{definition}

\begin{example}[Double dual]\label{ex:double-dual}
\index{natural transformation!double dual}
\index{double dual}
Let $\Vect_\K^{\mathrm{fd}}$ denote the category of finite-dimensional
$\K$-vector spaces.  Define $\eta\colon\Id\Rightarrow(-)^{**}$ by
\[
    \eta_V\colon V\longrightarrow V^{**}, \qquad
    v\longmapsto\bigl(\varphi\mapsto\varphi(v)\bigr).
\]
This is a natural isomorphism (each $\eta_V$ is an isomorphism by
dimension counting).  Crucially, the isomorphism $V\cong V^{**}$ is
\emph{natural} in~$V$, whereas the isomorphism $V\cong V^*$ (which
requires choosing a basis) is not.
\end{example}

\begin{example}[Determinant]\label{ex:determinant-nat-trans}
\index{natural transformation!determinant}
\index{determinant!as natural transformation}
Let $\Ring$ be the category of commutative rings.  Consider the functors
$\mathrm{GL}_n,\;(-)^\times\colon\Ring\to\Grp$ sending a
commutative ring~$R$ to the group of invertible $n\times n$ matrices over~$R$
and to the group of units~$R^\times$, respectively.
The determinant gives a natural transformation
\[
    \det\colon\mathrm{GL}_n\Longrightarrow(-)^\times.
\]
Naturality says: for every ring homomorphism $\varphi\colon R\to S$ and
every $M\in\mathrm{GL}_n(R)$,
$\varphi(\det M) = \det(\varphi(M))$.
\end{example}

\begin{definition}[Vertical composition]\label{def:vert-comp}
\index{natural transformation!vertical composition}
\index{vertical composition}
Let $\alpha\colon F\Rightarrow G$ and $\beta\colon G\Rightarrow H$ be
natural transformations between functors
$F,G,H\colon\mathcal{C}\to\mathcal{D}$.
Their \textbf{vertical composition}
$\beta\circ\alpha\colon F\Rightarrow H$ has components
\[
    (\beta\circ\alpha)_A = \beta_A\circ\alpha_A.
\]
\end{definition}

This is depicted schematically as:
\[
\begin{tikzcd}[column sep=huge]
    \mathcal{C}
    \arrow[r, bend left=50, "F"{name=F}]
    \arrow[r, "G"{name=G, description}]
    \arrow[r, bend right=50, "H"{name=H, below}]
    &
    \mathcal{D}
    \arrow[Rightarrow, from=F, to=G, "\alpha"]
    \arrow[Rightarrow, from=G, to=H, "\beta"]
\end{tikzcd}
\]

\begin{lemma}\label{lem:vert-comp-nat}
Vertical composition is associative, and the identity natural
transformation $\id_F\colon F\Rightarrow F$ (with components
$(\id_F)_A = \id_{F(A)}$) serves as identity.
\end{lemma}

\begin{proof}
Both statements follow immediately from the corresponding properties
of morphism composition in~$\mathcal{D}$.
\end{proof}

\begin{definition}[Horizontal composition]\label{def:horiz-comp}
\index{natural transformation!horizontal composition}
\index{horizontal composition}
\index{Godement product|see{horizontal composition}}
Let $\alpha\colon F\Rightarrow G$ between functors
$\mathcal{C}\to\mathcal{D}$ and
$\beta\colon H\Rightarrow K$ between functors
$\mathcal{D}\to\mathcal{E}$.
The \textbf{horizontal composition} (or \emph{Godement product})
$\beta*\alpha\colon H\circ F\Rightarrow K\circ G$ has components
\[
    (\beta*\alpha)_A = \beta_{G(A)}\circ H(\alpha_A)
                     = K(\alpha_A)\circ\beta_{F(A)}.
\]
The two expressions are equal by naturality of~$\beta$:
\[
\begin{tikzcd}
    HF(A) \arrow[r, "H(\alpha_A)"] \arrow[d, "\beta_{F(A)}"'] &
    HG(A) \arrow[d, "\beta_{G(A)}"] \\
    KF(A) \arrow[r, "K(\alpha_A)"'] &
    KG(A)
\end{tikzcd}
\]
\end{definition}


%% ============================================================
\section{Functor categories}\label{sec:functor-categories}
%% ============================================================

\begin{definition}[Functor category]\label{def:functor-category}
\index{functor category}
\index{category!functor}
Let $\mathcal{C}$ and $\mathcal{D}$ be categories with $\mathcal{C}$
small.  The \textbf{functor category}
$[\mathcal{C},\mathcal{D}]$ (also written
$\Fun(\mathcal{C},\mathcal{D})$ or $\mathcal{D}^{\mathcal{C}}$) is
the category whose:
\begin{itemize}
    \item objects are functors $F\colon\mathcal{C}\to\mathcal{D}$;
    \item morphisms are natural transformations;
    \item composition is vertical composition of natural
          transformations.
\end{itemize}
\end{definition}

\begin{remark}\label{rem:functor-cat-size}
\index{functor category!size}
If $\mathcal{C}$ is small and $\mathcal{D}$ is locally small, then
$[\mathcal{C},\mathcal{D}]$ is locally small.
\end{remark}

\begin{example}[Presheaf categories]\label{ex:presheaf}
\index{presheaf}
\index{presheaf category}
A particularly important special case: let $\mathcal{C}$ be a small
category.  The functor category
\[
    \widehat{\mathcal{C}} \;=\; [\mathcal{C}^{\op},\Set]
\]
is called the \textbf{category of presheaves} on~$\mathcal{C}$.  Its
objects are contravariant functors $\mathcal{C}\to\Set$, called
\emph{presheaves}.  This category has remarkable properties (it is
complete, cocomplete, and cartesian closed) that we shall explore later.
\end{example}

\begin{proposition}\label{prop:nat-iso-pointwise}
\index{natural isomorphism!pointwise criterion}
A natural transformation $\alpha\colon F\Rightarrow G$ is an isomorphism
in $[\mathcal{C},\mathcal{D}]$ if and only if each component
$\alpha_A$ is an isomorphism in~$\mathcal{D}$.
\end{proposition}

\begin{proof}
If $\alpha$ is a natural isomorphism, its inverse $\alpha^{-1}$ has
components $(\alpha^{-1})_A = (\alpha_A)^{-1}$.  One checks naturality
of $\alpha^{-1}$ by applying $(-\,)^{-1}$ to the naturality squares
of~$\alpha$.
Conversely, if each $\alpha_A$ is an isomorphism, define
$\beta_A = (\alpha_A)^{-1}$.  The naturality of~$\beta$ follows from
that of~$\alpha$: for $f\colon A\to B$,
\[
    \beta_B\circ G(f) = \beta_B\circ G(f)\circ\alpha_A\circ\beta_A
    = \beta_B\circ\alpha_B\circ F(f)\circ\beta_A
    = F(f)\circ\beta_A. \qedhere
\]
\end{proof}

\begin{definition}[Equivalence of categories]\label{def:equivalence}
\index{equivalence of categories}
\index{category!equivalence}
An \textbf{equivalence of categories} between $\mathcal{C}$ and
$\mathcal{D}$ consists of functors
$F\colon\mathcal{C}\to\mathcal{D}$ and
$G\colon\mathcal{D}\to\mathcal{C}$ together with natural isomorphisms
\[
    \eta\colon\Id_\mathcal{C}\xRightarrow{\;\sim\;}G\circ F
    \qquad\text{and}\qquad
    \varepsilon\colon F\circ G\xRightarrow{\;\sim\;}\Id_\mathcal{D}.
\]
We write $\mathcal{C}\simeq\mathcal{D}$ and say $\mathcal{C}$ and
$\mathcal{D}$ are \emph{equivalent}.
\end{definition}

\begin{theorem}\label{thm:equiv-ff-eso}
\index{equivalence of categories!characterisation}
A functor $F\colon\mathcal{C}\to\mathcal{D}$ is part of an equivalence
of categories if and only if $F$ is fully faithful and essentially
surjective.
\end{theorem}

\begin{proof}
($\Rightarrow$) Suppose $F$ is an equivalence with quasi-inverse~$G$ and
natural isomorphisms $\eta\colon\Id\xRightarrow{\sim} GF$ and
$\varepsilon\colon FG\xRightarrow{\sim}\Id$.

\emph{Essentially surjective:} For $D\in\mathcal{D}$,
$F(G(D))\cong D$ via $\varepsilon_D$.

\emph{Faithful:} Suppose $F(f)=F(g)$ for $f,g\colon A\to B$.
Then $GF(f)=GF(g)$.  Naturality of~$\eta$ gives
$\eta_B\circ f = GF(f)\circ\eta_A = GF(g)\circ\eta_A = \eta_B\circ g$,
and since $\eta_B$ is an isomorphism, $f=g$.

\emph{Full:} Let $h\colon F(A)\to F(B)$.  Consider
$f = \eta_B^{-1}\circ G(h)\circ\eta_A\colon A\to B$.
Then $F(f) = \varepsilon_{F(B)}\circ FG(h)\circ(\varepsilon_{F(A)})^{-1}$.
By naturality of~$\varepsilon$,
$\varepsilon_{F(B)}\circ FG(h) = h\circ\varepsilon_{F(A)}$,
so $F(f) = h$.

($\Leftarrow$)  Assume $F$ is fully faithful and essentially surjective.
For each $D\in\mathcal{D}$, choose $G(D)\in\mathcal{C}$ with an
isomorphism $\varepsilon_D\colon FG(D)\xrightarrow{\sim} D$.
For $h\colon D\to D'$, define $G(h)$ to be the unique morphism
(existing by full faithfulness) satisfying
$F(G(h)) = \varepsilon_{D'}^{-1}\circ h\circ\varepsilon_D$.
One verifies $G$ is a functor and that $\varepsilon$ and the induced
$\eta$ are natural isomorphisms.
\end{proof}


%% ============================================================
\section{Exercises for Chapter~1}\label{sec:exercises-ch1}
%% ============================================================

\begin{exercise}\label{exo:cat-rels}
\index{category!of relations}
Define a category $\mathbf{Rel}$ whose objects are sets and whose
morphisms $A\to B$ are relations $R\subseteq A\times B$.
What is composition?  What are the identity morphisms?
Show that $\mathbf{Rel}$ is indeed a category.
\end{exercise}

\begin{exercise}\label{exo:groupoid}
\index{groupoid}
A \emph{groupoid} is a category in which every morphism is an
isomorphism.  Show that the following are groupoids:
\begin{enumerate}[label=(\roman*)]
    \item A group, viewed as a one-object category.
    \item The \emph{fundamental groupoid} $\Pi_1(X)$ of a topological
          space~$X$: objects are points of~$X$, morphisms $x\to y$
          are homotopy classes of paths from~$x$ to~$y$.
    \item The \emph{core} $\mathrm{core}(\mathcal{C})$ of any category
          $\mathcal{C}$: same objects, but only isomorphisms.
\end{enumerate}
\end{exercise}

\begin{exercise}\label{exo:slice-category}
\index{slice category}
\index{category!slice}
Let $\mathcal{C}$ be a category and $A\in\mathcal{C}$ an object.
The \emph{slice category} (or \emph{over category}) $\mathcal{C}/A$
has:
\begin{itemize}
    \item Objects: morphisms $f\colon X\to A$ in~$\mathcal{C}$.
    \item Morphisms: from $(f\colon X\to A)$ to $(g\colon Y\to A)$,
          a morphism $h\colon X\to Y$ in~$\mathcal{C}$ such that
          $g\circ h = f$:
\[
\begin{tikzcd}
    X \arrow[rr, "h"] \arrow[dr, "f"'] && Y \arrow[dl, "g"] \\
    & A &
\end{tikzcd}
\]
\end{itemize}
Verify that $\mathcal{C}/A$ is a category.  What is the terminal object?
Describe $\Set/\{0,1\}$ explicitly.
\end{exercise}

\begin{exercise}\label{exo:comma-category}
\index{comma category}
Let $F\colon\mathcal{A}\to\mathcal{C}$ and
$G\colon\mathcal{B}\to\mathcal{C}$ be functors.
Define the \emph{comma category} $(F\downarrow G)$ and show that
slice categories and coslice categories are special cases.
\end{exercise}

\begin{exercise}\label{exo:mono-epi-exercises}
\begin{enumerate}[label=(\roman*)]
    \item Show that in any category, the composition of two
          monomorphisms is a monomorphism.
    \item Show that in $\Set$, a morphism is an epimorphism if and
          only if it is surjective.
    \item Find an example of a morphism in a concrete category that is
          both monic and epic but not an isomorphism
          (beyond the example of $\Z\hookrightarrow\Q$ in $\Ring$).
\end{enumerate}
\end{exercise}

\begin{exercise}\label{exo:endomorphism-monoid}
\index{endomorphism!monoid}
For any object $A$ in a category~$\mathcal{C}$, show that
$\End_\mathcal{C}(A)=\Hom_\mathcal{C}(A,A)$ forms a monoid under
composition.  When is this monoid a group?
\end{exercise}

\begin{exercise}\label{exo:functor-composition}
\index{functor!composition}
Let $F\colon\mathcal{A}\to\mathcal{B}$ and
$G\colon\mathcal{B}\to\mathcal{C}$ be functors.
Show that the composition $G\circ F\colon\mathcal{A}\to\mathcal{C}$
is again a functor.  Verify that functor composition is associative
and that the identity functor $\Id_\mathcal{C}$ is a two-sided identity.
\end{exercise}

\begin{exercise}\label{exo:nat-trans-id-functor}
\index{natural transformation!from identity functor}
Let $F\colon\mathcal{C}\to\mathcal{C}$ be a functor.
Show that the following are equivalent:
\begin{enumerate}[label=(\roman*)]
    \item $\alpha\colon\Id_\mathcal{C}\Rightarrow F$ is a natural
          transformation;
    \item For every morphism $f\colon A\to B$, one has
          $F(f)\circ\alpha_A = \alpha_B\circ f$.
\end{enumerate}
\end{exercise}

\begin{exercise}\label{exo:interchange-law}
\index{interchange law}
\index{natural transformation!interchange law}
Prove the \emph{interchange law} for horizontal and vertical
composition of natural transformations.  Specifically, given
\[
\begin{tikzcd}[column sep=large]
    \mathcal{A}
    \arrow[r, bend left=40, "F_1"{name=F1}]
    \arrow[r, bend right=40, "F_2"{name=F2, below}]
    &
    \mathcal{B}
    \arrow[r, bend left=40, "G_1"{name=G1}]
    \arrow[r, bend right=40, "G_2"{name=G2, below}]
    &
    \mathcal{C}
    \arrow[Rightarrow, from=F1, to=F2, "\alpha"]
    \arrow[Rightarrow, from=G1, to=G2, "\beta"]
\end{tikzcd}
\]
and natural transformations $\alpha'\colon F_2\Rightarrow F_3$,
$\beta'\colon G_2\Rightarrow G_3$, show that
\[
    (\beta'\circ\beta)*(\alpha'\circ\alpha)
    = (\beta'*\alpha')\circ(\beta*\alpha).
\]
\end{exercise}

\begin{exercise}\label{exo:skeleton}
\index{skeleton of a category}
A category~$\mathcal{C}$ is \emph{skeletal} if isomorphic objects are
equal.  A \emph{skeleton} of~$\mathcal{C}$ is a full subcategory that
is skeletal and contains exactly one object from each isomorphism class.
Show that every category has a skeleton and that inclusion of the
skeleton is an equivalence of categories.
\end{exercise}

\begin{exercise}\label{exo:nat-trans-precomposition}
Let $F,G\colon\mathcal{C}\to\mathcal{D}$ be functors,
$\alpha\colon F\Rightarrow G$ a natural transformation, and
$H\colon\mathcal{B}\to\mathcal{C}$ a functor.
Define the \emph{whiskering} $\alpha H\colon FH\Rightarrow GH$ by
$(\alpha H)_B = \alpha_{H(B)}$.
Show that $\alpha H$ is a natural transformation.
Similarly define whiskering on the other side.
\end{exercise}

\index{functor|)}
\index{natural transformation|)}
\index{category|)}


%%%%%%%%%%%%%%%%%%%%%%%%%%%%%%%%%%%%%%%%%%%%%%%%%%%%%%%%%%%%
%%       CHAPTER 2: DUALITY AND THE DUAL PRINCIPLE        %%
%%%%%%%%%%%%%%%%%%%%%%%%%%%%%%%%%%%%%%%%%%%%%%%%%%%%%%%%%%%%

\chapter{Duality and the Dual Principle}\label{ch:duality}
\minitoc

\vspace{1em}

One of the most powerful features of category theory is the
\emph{duality principle}: every categorical statement has a dual,
obtained by reversing all morphisms.  This single observation doubles the
theorems one can prove ``for free''.  In this chapter we make this
precise by introducing opposite categories, formulating the duality
principle, and examining several concrete dualities.  We also study the
two-variable hom-functor, which will be essential in later chapters.

\index{duality|(}


%% ============================================================
\section{The opposite category}\label{sec:opposite-cat}
%% ============================================================

\begin{definition}[Opposite category]\label{def:opposite-cat}
\index{opposite category}
\index{category!opposite}
\index{dual category|see{opposite category}}
Let $\mathcal{C}$ be a category.  The \textbf{opposite category} (or
\emph{dual category}) $\mathcal{C}^{\op}$ is defined by:
\begin{enumerate}[label=(\roman*)]
    \item $\Ob(\mathcal{C}^{\op}) = \Ob(\mathcal{C})$.
    \item For objects $A,B$:
          $\Hom_{\mathcal{C}^{\op}}(A,B) = \Hom_\mathcal{C}(B,A)$.
    \item Composition in $\mathcal{C}^{\op}$:
          given $f^{\op}\colon A\to B$ and $g^{\op}\colon B\to C$
          in~$\mathcal{C}^{\op}$ (corresponding to $f\colon B\to A$ and
          $g\colon C\to B$ in~$\mathcal{C}$), we set
          $g^{\op}\circ f^{\op} = (f\circ g)^{\op}$.
    \item The identity on~$A$ in $\mathcal{C}^{\op}$ is $(\id_A)^{\op}
          = \id_A$.
\end{enumerate}
\end{definition}

\begin{proposition}\label{prop:op-involution}
\index{opposite category!involution}
For any category~$\mathcal{C}$, $(\mathcal{C}^{\op})^{\op}=\mathcal{C}$.
\end{proposition}

\begin{proof}
The objects coincide.  For morphisms,
$\Hom_{(\mathcal{C}^{\op})^{\op}}(A,B)
= \Hom_{\mathcal{C}^{\op}}(B,A)
= \Hom_\mathcal{C}(A,B)$.
Composition: $(f^{\op})^{\op}\circ(g^{\op})^{\op}
= ((g^{\op}\circ f^{\op})^{\op})
= ((f\circ g)^{\op})^{\op}
= f\circ g$.
Everything reduces to the original data of~$\mathcal{C}$.
\end{proof}

\begin{example}\label{ex:op-examples}
\begin{enumerate}[label=(\roman*)]
    \item If $(P,\le)$ is a poset viewed as a category, then
          $P^{\op}$ is the poset $(P,\ge)$.
          \index{poset!opposite}
    \item If $G$ is a group viewed as a one-object category, then
          $G^{\op}$ is the \emph{opposite group}: same elements,
          with multiplication $a\cdot^{\op}b = b\cdot a$.
          For abelian groups, $G^{\op}\cong G$.
          \index{group!opposite}
    \item $\Set^{\op}$ is a perfectly valid category, but it is not
          equivalent to any ``familiar'' category of structured sets
          (it is, however, equivalent to the category of complete
          atomic Boolean algebras, by a non-trivial theorem).
          \index{category!Set op@$\Set^{\op}$}
\end{enumerate}
\end{example}


%% ============================================================
\section{The duality principle}\label{sec:duality-principle}
%% ============================================================

\begin{definition}[Dual statement]\label{def:dual-statement}
\index{dual statement}
\index{duality principle}
Let $\Sigma$ be a statement formulated purely in the language of
category theory (objects, morphisms, composition, identities).
The \textbf{dual statement}~$\Sigma^{\op}$ is obtained from~$\Sigma$
by:
\begin{itemize}
    \item reversing the direction of every morphism;
    \item replacing each composition $g\circ f$ by $f\circ g$;
    \item interchanging domain and codomain.
\end{itemize}
\end{definition}

\begin{theorem}[Duality principle]\label{thm:duality-principle}
\index{duality principle!statement}
If a statement~$\Sigma$ holds in every category, then so does its
dual~$\Sigma^{\op}$.
More generally, if $\Sigma$ holds in a category~$\mathcal{C}$, then
$\Sigma^{\op}$ holds in~$\mathcal{C}^{\op}$.
\end{theorem}

\begin{proof}
A statement about~$\mathcal{C}^{\op}$ is, by definition, a statement
about~$\mathcal{C}$ with all arrows reversed.
Since $\mathcal{C}^{\op}$ is a legitimate category, any theorem valid
for all categories applies to it.
Unwinding the definitions in $\mathcal{C}^{\op}$ recovers the dual
statement in~$\mathcal{C}$.
\end{proof}

\begin{remark}\label{rem:duality-examples}
\index{duality principle!examples}
The duality principle immediately gives us dual pairs of concepts:
\begin{center}
\begin{tabular}{ll}
\toprule
\textbf{Concept} & \textbf{Dual concept} \\
\midrule
monomorphism & epimorphism \\
initial object & terminal object \\
product & coproduct \\
limit & colimit \\
kernel & cokernel \\
left adjoint & right adjoint \\
\bottomrule
\end{tabular}
\end{center}
Once we prove a theorem about monomorphisms, the dual theorem about
epimorphisms follows automatically.
\end{remark}

\begin{example}\label{ex:duality-mono-epi}
\index{monomorphism!dual of epimorphism}
\index{epimorphism!dual of monomorphism}
We showed in 
ef{lem:mono-epi-composition} that if $g\circ f$ is monic
then $f$ is monic.  By duality: if $g\circ f$ is epic, then $g$ is
epic.  No additional proof is needed---it follows from the proof of the
original statement applied to $\mathcal{C}^{\op}$.
\end{example}


%% ============================================================
\section{Functors and duality}\label{sec:functors-duality}
%% ============================================================

\begin{proposition}\label{prop:contravariant-op}
\index{functor!contravariant!via opposite category}
A contravariant functor $F\colon\mathcal{C}\to\mathcal{D}$ is the same
thing as a covariant functor $\mathcal{C}^{\op}\to\mathcal{D}$, and
also the same as a covariant functor $\mathcal{C}\to\mathcal{D}^{\op}$.
\end{proposition}

\begin{proof}
A covariant functor $\mathcal{C}^{\op}\to\mathcal{D}$ assigns to each
morphism $f^{\op}\colon B\to A$ in $\mathcal{C}^{\op}$ (i.e.\
$f\colon A\to B$ in $\mathcal{C}$) a morphism
$F(f^{\op})\colon F(B)\to F(A)$ in~$\mathcal{D}$,
which is precisely a contravariant assignment.
The preservation of composition and identities translates
directly.  The argument for $\mathcal{C}\to\mathcal{D}^{\op}$ is
analogous.
\end{proof}

\begin{remark}\label{rem:op-functor}
\index{opposite functor}
Every functor $F\colon\mathcal{C}\to\mathcal{D}$ induces a functor
$F^{\op}\colon\mathcal{C}^{\op}\to\mathcal{D}^{\op}$ defined by
$F^{\op}(A)=F(A)$ on objects and $F^{\op}(f^{\op})=(F(f))^{\op}$ on
morphisms.  In fact, $(-)^{\op}$ is a functor $\Cat\to\Cat$.
\end{remark}


%% ============================================================
\section{Concrete dualities}\label{sec:concrete-dualities}
%% ============================================================

\index{duality!concrete|(}

The abstract duality principle tells us about dual \emph{concepts}, but
there also exist concrete \emph{equivalences}
$\mathcal{C}^{\op}\simeq\mathcal{D}$ for specific categories.  These are
called \emph{concrete dualities} or \emph{dualities of categories}.

\begin{example}[Stone duality]\label{ex:stone-duality}
\index{Stone duality}
\index{Boolean algebra}
\index{Stone space}
The category $\mathbf{Stone}$ of Stone spaces (compact, Hausdorff,
totally disconnected topological spaces) and continuous maps is
equivalent to $\mathbf{Bool}^{\op}$, where $\mathbf{Bool}$ is the
category of Boolean algebras and their homomorphisms:
\[
    \mathbf{Stone} \;\simeq\; \mathbf{Bool}^{\op}.
\]
The equivalence is implemented by the functor sending a Stone space to
its Boolean algebra of clopen subsets, and conversely sending a
Boolean algebra to its space of ultrafilters.
\end{example}

\begin{example}[Pontryagin duality]\label{ex:pontryagin-duality}
\index{Pontryagin duality}
\index{locally compact abelian group}
Let $\mathbf{LCA}$ be the category of locally compact abelian groups
and continuous homomorphisms.  Pontryagin duality provides a
contravariant equivalence
\[
    (-)^\wedge\colon\mathbf{LCA}^{\op}
    \xrightarrow{\;\sim\;}\mathbf{LCA},
\]
where $G^\wedge = \Hom_{\mathrm{cts}}(G,\R/\Z)$ is the
\emph{Pontryagin dual}.
The natural map $G\to G^{\wedge\wedge}$, $g\mapsto(\chi\mapsto\chi(g))$,
is an isomorphism.
\end{example}

\begin{example}[Gelfand duality]\label{ex:gelfand-duality}
\index{Gelfand duality}
\index{C*-algebra@$C^*$-algebra}
There is an equivalence
\[
    \mathbf{CHaus}^{\op} \;\simeq\; \mathbf{cC^*Alg}_1,
\]
where $\mathbf{CHaus}$ is the category of compact Hausdorff spaces and
$\mathbf{cC^*Alg}_1$ is the category of commutative unital
$C^*$-algebras.  The functor sends $X\mapsto C(X,\C)$
(continuous functions) and the quasi-inverse sends a $C^*$-algebra to
its maximal ideal space (Gelfand spectrum).
\end{example}

\begin{example}[Finite-dimensional vector space duality]\label{ex:fd-vect-duality}
\index{duality!finite-dimensional vector spaces}
\index{vector space!duality}
The dual space functor $(-)^*\colon\Vect_\K^{\mathrm{fd}}\to
(\Vect_\K^{\mathrm{fd}})^{\op}$ is an equivalence of categories.
Combined with $(\Vect_\K^{\mathrm{fd}})^{\op}\simeq
\Vect_\K^{\mathrm{fd}}$
(via the double dual), we see that finite-dimensional vector spaces
are self-dual.
\end{example}

\index{duality!concrete|)}


%% ============================================================
\section{The hom-functor in both variables}\label{sec:hom-bifunctor}
%% ============================================================

\begin{definition}[Bifunctor]\label{def:bifunctor}
\index{bifunctor}
\index{functor!of two variables}
Let $\mathcal{C}$, $\mathcal{D}$, $\mathcal{E}$ be categories.
A \textbf{bifunctor} is a functor
$F\colon\mathcal{C}\times\mathcal{D}\to\mathcal{E}$.
\end{definition}

\begin{remark}\label{rem:bifunctor-partial}
A bifunctor $F\colon\mathcal{C}\times\mathcal{D}\to\mathcal{E}$
determines, for each $C\in\mathcal{C}$, a functor
$F(C,-)\colon\mathcal{D}\to\mathcal{E}$ and, for each
$D\in\mathcal{D}$, a functor
$F(-,D)\colon\mathcal{C}\to\mathcal{E}$.
Conversely, a family of functors in each variable that is
``jointly functorial'' determines a bifunctor.
\end{remark}

\begin{proposition}\label{prop:hom-bifunctor}
\index{hom-functor!bifunctor}
\index{hom-bifunctor}
For any locally small category~$\mathcal{C}$, the assignment
\[
    \Hom_\mathcal{C}(-,-)\colon\mathcal{C}^{\op}\times\mathcal{C}
    \longrightarrow\Set
\]
defined on objects by $(A,B)\mapsto\Hom_\mathcal{C}(A,B)$ and on
morphisms by
\[
    (f\colon A'\to A,\; g\colon B\to B')\;\longmapsto\;
    \bigl(h\mapsto g\circ h\circ f\bigr)
    \colon\Hom(A,B)\to\Hom(A',B'),
\]
is a bifunctor (i.e.\ a functor from the product category
$\mathcal{C}^{\op}\times\mathcal{C}$ to $\Set$).
\end{proposition}

\begin{proof}
We verify the functor axioms.  On identities:
\[
    \Hom(\id_A,\id_B)(h) = \id_B\circ h\circ\id_A = h,
\]
so $\Hom(\id_A,\id_B) = \id_{\Hom(A,B)}$.

On composition: let $f'\colon A''\to A'$, $f\colon A'\to A$,
$g\colon B\to B'$, $g'\colon B'\to B''$.
We need
$\Hom(f\circ f',\,g'\circ g) = \Hom(f',g')\circ\Hom(f,g)$.
The left side maps $h\mapsto (g'\circ g)\circ h\circ(f\circ f')$.
The right side maps
$h\mapsto g\circ h\circ f\mapsto g'\circ(g\circ h\circ f)\circ f'$,
which is the same by associativity.
\end{proof}

\begin{remark}\label{rem:hom-functor-partial-functors}
\index{hom-functor!covariant}
\index{hom-functor!contravariant}
The bifunctor $\Hom(-,-)$ encodes the two ``partial'' hom-functors
from 
ef{ex:hom-functors}:
\begin{itemize}
    \item Fixing the first argument: $\Hom(A,-)$ is covariant.
    \item Fixing the second argument: $\Hom(-,B)$ is contravariant
          (i.e.\ covariant from $\mathcal{C}^{\op}$).
\end{itemize}
The two partial functors determine the bifunctor (they agree on pairs
of identities and satisfy a compatibility condition that is automatic
from the bifunctor structure).
\end{remark}

\begin{proposition}\label{prop:hom-preserves-limits}
\index{hom-functor!preserves limits}
For any object $A$ in a locally small category~$\mathcal{C}$, the
covariant hom-functor $\Hom(A,-)\colon\mathcal{C}\to\Set$ preserves
all limits that exist in~$\mathcal{C}$.
\end{proposition}

\begin{proof}
This is a consequence of the universal property of limits and will be
proved rigorously in 
ef{ch:limits-colimits}.  The key idea is that a
natural bijection $\Hom(A,\lim F)\cong\lim\Hom(A,F-)$ follows directly
from the defining universal property of the limit.
\end{proof}

\begin{definition}[Representable functor (preview)]\label{def:representable-preview}
\index{representable functor!preview}
\index{functor!representable}
A functor $F\colon\mathcal{C}\to\Set$ is \textbf{representable} if it
is naturally isomorphic to $\Hom_\mathcal{C}(A,-)$ for some object
$A\in\mathcal{C}$.  The object~$A$ is called a \emph{representing
object}.

A contravariant functor $F\colon\mathcal{C}^{\op}\to\Set$ is
representable if $F\cong\Hom_\mathcal{C}(-,B)$ for some~$B$.

Representable functors are central to category theory; we shall study
them in depth in connection with the Yoneda lemma.
\end{definition}


%% ============================================================
\section{Exercises for Chapter~2}\label{sec:exercises-ch2}
%% ============================================================

\begin{exercise}\label{exo:op-op}
Let $\mathcal{C}$ be a category.
\begin{enumerate}[label=(\roman*)]
    \item Verify carefully that $\mathcal{C}^{\op}$ as defined in
          
ef{def:opposite-cat} satisfies the axioms of a category.
    \item Show that $(\mathcal{C}^{\op})^{\op}=\mathcal{C}$ (equality,
          not merely isomorphism).
\end{enumerate}
\end{exercise}

\begin{exercise}\label{exo:initial-terminal-dual}
\index{initial object!dual of terminal}
Show that $A$ is an initial object in $\mathcal{C}$ if and only if
$A$ is a terminal object in $\mathcal{C}^{\op}$.
Deduce the uniqueness (up to unique isomorphism) of terminal objects
from the corresponding result for initial objects via duality.
\end{exercise}

\begin{exercise}\label{exo:mono-epi-dual}
Show that $f$ is a monomorphism in $\mathcal{C}$ if and only if
$f^{\op}$ is an epimorphism in $\mathcal{C}^{\op}$.
\end{exercise}

\begin{exercise}\label{exo:op-functor-category}
\index{functor category!opposite}
Let $\mathcal{C}$ be small and $\mathcal{D}$ any category.
Construct a canonical isomorphism of categories
\[
    [\mathcal{C},\mathcal{D}]^{\op}
    \;\cong\;
    [\mathcal{C},\mathcal{D}^{\op}].
\]
(Hint: what happens to naturality squares when you reverse arrows in the
target?)
\end{exercise}

\begin{exercise}\label{exo:self-dual-category}
\index{self-dual category}
A category~$\mathcal{C}$ is called \emph{self-dual} if
$\mathcal{C}\simeq\mathcal{C}^{\op}$.
\begin{enumerate}[label=(\roman*)]
    \item Show that $\Vect_\K^{\mathrm{fd}}$ is self-dual.
    \item Show that any groupoid is self-dual.
    \item Is $\Set$ self-dual?  Justify.
\end{enumerate}
\end{exercise}

\begin{exercise}\label{exo:product-category}
\index{product category}
Let $\mathcal{C}$ and $\mathcal{D}$ be categories.
Define the product category $\mathcal{C}\times\mathcal{D}$ and show that
\[
    (\mathcal{C}\times\mathcal{D})^{\op}
    \;\cong\;
    \mathcal{C}^{\op}\times\mathcal{D}^{\op}.
\]
\end{exercise}

\begin{exercise}\label{exo:hom-bifunctor-verification}
\index{hom-bifunctor!detailed verification}
Let $\mathcal{C}$ be a locally small category.
Give a detailed verification that
$\Hom_\mathcal{C}(-,-)\colon\mathcal{C}^{\op}\times\mathcal{C}\to\Set$
is a functor by checking the functor axioms (preservation of identities
and composition) explicitly.
\end{exercise}

\begin{exercise}\label{exo:contravariant-nat-trans}
\index{natural transformation!between contravariant functors}
Let $F,G\colon\mathcal{C}^{\op}\to\Set$ be contravariant functors.
Spell out explicitly what a natural transformation
$\alpha\colon F\Rightarrow G$ looks like: what are the components, and
what does the naturality condition say?
Draw the naturality square and compare it with the covariant case.
\end{exercise}

\begin{exercise}\label{exo:op-preserves-equivalence}
\index{equivalence of categories!and opposite}
Show that if $\mathcal{C}\simeq\mathcal{D}$, then
$\mathcal{C}^{\op}\simeq\mathcal{D}^{\op}$.
\end{exercise}

\begin{exercise}\label{exo:concrete-duality-exploration}
\index{duality!in lattices}
Let $L$ be a bounded lattice, viewed as a category (poset).
Show that $L^{\op}$ is again a bounded lattice, with joins and meets
interchanged.  Describe the dual of the lattice of open sets of a
topological space~$X$.
\end{exercise}

\begin{exercise}\label{exo:op-slice}
\index{slice category!opposite}
\index{coslice category}
Let $\mathcal{C}$ be a category and $A\in\mathcal{C}$ an object.
The \emph{coslice category} (or \emph{under category}) $A/\mathcal{C}$
has objects $f\colon A\to X$ and morphisms the evident commutative
triangles.
Show that $(A/\mathcal{C})^{\op}\cong\mathcal{C}^{\op}/A$.
\end{exercise}

\begin{exercise}\label{exo:representable-dual}
\index{representable functor!and duality}
Let $\mathcal{C}$ be a locally small category and $A\in\mathcal{C}$.
Show that the contravariant hom-functor $\Hom(-,A)$ is representable
as a functor $\mathcal{C}^{\op}\to\Set$, with representing object~$A$
(viewed as an object of $\mathcal{C}^{\op}$).
\end{exercise}

\index{duality|)}
